% ===== Hallucinations and delusions: "Reality is overrated anyway" (~17 min) =====
\partaccent{brainpink}
\topicopener{speakinghead}{Հալյուցինացիաներ}{Իրականությունը, միևնույն է, գերագնահատված է}{Հալյուցինացիաներ, զառանցանքներ, և ինչ է անում ուղեղը դրանք առաջացնելու համար}{thinking}

\begin{frame}{Փսիխոզ}
  \relax{\small \textbf{Փսիխոզ} (psychosis)՝ վիճակ, երբ մարդը դժվարանում է տարբերել իրականն անիրականից}\par\medskip
  \begin{columns}[T]
    \begin{column}{0.47\textwidth}
      \begin{card}[colback=lavender!40!cardwhite]
        \emoji[1.3em]{ear}\ {\blackfont Հալյուցինացիա}\par\smallskip
        Ընկալել այն, ինչ \textbf{չկա}. տեսնել, լսել, զգալ
      \end{card}
    \end{column}
    \begin{column}{0.47\textwidth}
      \onslide<2->{%
      \begin{card}[colback=peach!40!cardwhite]
        \emoji[1.3em]{thought}\ {\blackfont Զառանցանք} (delusion)\par\smallskip
        Անսասան հավատալ մի բանի, որն ակնհայտորեն \textbf{ճիշտ չէ}
      \end{card}}
    \end{column}
  \end{columns}
  \vspace{3mm}
  \onslide<3->{%
  \begin{card}[colback=butter!45!cream]
    \centering Իրականությունը կորցնելու միտքն ինքնին վախեցնող է։ Իսկ ինչպե՞ս ապրել դրա հետ
  \end{card}}
\end{frame}

\begin{frame}{Իրականության համակարգը շատ խոցելի է}
  \relax{\small Հալյուցինացիաներ ու զառանցանքներ կարող են առաջացնել.}\par\medskip
  \chip[lavender]{raincloud}{դեպրեսիա}\ \chip[lavender]{battery}{սթրես, պոռթկում}\ \chip[lavender]{monkey}{թմրանյութեր, ալկոհոլ}\par\smallskip
  {\scriptsize\color{plum!70!cream} այսօրվա բոլոր թեմաները, ու նաև.}\par\smallskip
  \onslide<2->{%
  \chip{}{դեմենցիա}\ \chip{}{Պարկինսոն}\ \chip{}{բիպոլյար խանգարում}\ \chip{zzz}{քնի պակաս}\ \chip{}{ուղեղի ուռուցք}\par\smallskip
  \chip{}{ՄԻԱՎ}\ \chip{}{սիֆիլիս}\ \chip{}{Լայմի հիվանդություն}\ \chip{}{ցրված սկլերոզ}\ \chip{candy}{ցածր շաքար}\par\smallskip
  \chip{}{կանեփ}\ \chip{}{ամֆետամիններ}\ \chip{}{կետամին}\ \chip{}{կոկաին}\par\medskip}
  \onslide<3->{%
  \begin{card}[colback=brainpink!25!cardwhite]
    Որոշ վիճակներ այնքան են կապված փսիխոզի հետ, որ կոչվում են \textbf{փսիխոտիկ խանգարումներ}։ Ամենահայտնին՝ \textbf{շիզոֆրենիան}
  \end{card}}
\end{frame}

\begin{frame}{Առասպել, թե՞ փաստ}
  \mythorfact{Շիզոֆրենիան «երկատված անհատականություն» է}{ԱՌԱՍՊԵԼ}{pinkdark}{%
    \begin{card}[colback=brainpink!25!cardwhite]
      Հունարեն «schizo»՝ «երկատել»։ Բայց երկատվածությունը ոչ թե անհատականությունների, այլ \textbf{մարդու ու իրականության} միջև է։
    \end{card}}
\end{frame}

\begin{frame}{«Ձայներ»}
  \relax{\small Ամենահաճախը \textbf{լսողական} հալյուցինացիաներն են. «ձայներ լսել»։ Կան նաև հպման, համի, հոտի հալյուցինացիաներ}\par\medskip
  \begin{columns}[T]
    \begin{column}{0.31\textwidth}
      \begin{card}[colback=lavender!40!cardwhite]
        {\blackfont 1-ին դեմք}\par\smallskip
        Սեփական մտքերդ «լսում ես», ասես ուրիշն է ասում
      \end{card}
    \end{column}
    \begin{column}{0.31\textwidth}
      \onslide<2->{%
      \begin{card}[colback=peach!40!cardwhite]
        {\blackfont 2-րդ դեմք}\par\smallskip
        Ձայնը խոսում է \textbf{քեզ հետ}. «Դու...»
      \end{card}}
    \end{column}
    \begin{column}{0.31\textwidth}
      \onslide<3->{%
      \begin{card}[colback=brainpink!30!cardwhite]
        {\blackfont 3-րդ դեմք}\par\smallskip
        Ձայները խոսում են \textbf{քո մասին}, մեկնաբանում են, թե ինչ ես անում
      \end{card}}
    \end{column}
  \end{columns}
  \vspace{2mm}
  \onslide<4->{\small Ձայները կարող են լինել տղամարդու կամ կնոջ, ծանոթ կամ անծանոթ, բարի կամ քննադատող։ Հաճախ՝ \textbf{քննադատող}։\par}
\end{frame}

\begin{frame}{Ինչպե՞ս ուսումնասիրել հալյուցինացիաները}
  \begin{columns}[c]
    \begin{column}{0.5\textwidth}
      \begin{itemize}
        \item Լաբորատորիայում մարդիկ չեն կարող հալյուցինացիա ունենալ «պատվերով»
        \item Եթե կարողանային դրանք միացնել-անջատել, խնդիր էլ չէր լինի
        \item Ուսումնասիրում են հիմնականում շիզոֆրենիայի համառ, լսողական հալյուցինացիաները
      \end{itemize}
    \end{column}
    \begin{column}{0.46\textwidth}
      \onslide<2->{%
      \begin{scienceupdate}
        Գրքում՝ համառ, քննադատող «3-րդ դեմքի» ձայները շիզոֆրենիայի հուսալի նշան են։ DSM-5-ը (2013) այլևս առանձին կշիռ չի տալիս դրանց. նման ձայներ լինում են նաև այլ վիճակներում։
      \end{scienceupdate}}
    \end{column}
  \end{columns}
  \only<2->{\source{Tandon et al., Schizophr. Res., 2013}}
\end{frame}

\begin{frame}{Ներքին ու արտաքին}
  \begin{columns}[c]
    \begin{column}{0.48\textwidth}
      \begin{itemize}
        \item Ուղեղն անընդհատ «շաղակրատում» է. մտածում, անհանգստանում, երազում
        \item Սովորաբար այն լավ տարբերում է \textbf{ներքին} ակտիվությունը \textbf{արտաքինից} (զգայարաններից)
        \item<3-> Գլխավոր վարկածը. հալյուցինացիայի ժամանակ այդ տարբերակումը խախտվում է
      \end{itemize}
    \end{column}
    \begin{column}{0.48\textwidth}
      \centering
      \begin{tikzpicture}[fold/.style={draw=plum, line width=0.8pt, rounded corners=1mm, minimum width=2.4cm, minimum height=1.25cm,
          font=\small, align=center, text=plum}]
        \only<1-2>{%
        \node[fold, fill=mint!45!cardwhite] (s) at (-1.45,0) {\emoji{memo}\\ ուղարկված\\ (ներքին)};
        \node[fold, fill=lavender!50!cardwhite] (r) at (1.45,0) {\emoji{loveletter}\\ ստացված\\ (արտաքին)};}
        \only<3->{%
        \node[fold, fill=brainpink!45!cardwhite, minimum width=4.4cm] (m) at (0,0) {\emoji{memo}\emoji{loveletter}\emoji{memo}\\ ամեն ինչ\\ \textbf{մեկ թղթապանակում}};}
      \end{tikzpicture}\par\medskip
      \onslide<2->{\alt<3->{{\small Թղթապանակները \textbf{խառնվել} են}}{{\small Ինչպես նամակները՝ առանձին թղթապանակներում}}\par}
      \onslide<4->{%
      \begin{card}[colback=butter!45!cream]
        Հիշո՞ւմ ես կարտոֆիլն ու խնձորը առաջին վեբինարից։ Ուղեղը հեշտությամբ է շփոթում աղբյուրները
      \end{card}}
    \end{column}
  \end{columns}
\end{frame}

\begin{frame}{Ինչու չենք կարողանում ինքներս մեզ խուտուտ տալ}
  \begin{columns}[c]
    \begin{column}{0.5\textwidth}
      \begin{tryit}
        Փորձիր հիմա ինքդ քեզ խուտուտ տալ ափիդ կամ կողքիդ։ Ստացվու՞մ է։ Գրիր չաթում։
      \end{tryit}
      \onslide<2->{%
      \begin{itemize}
        \item Ուղեղը \textbf{նախապես գիտի}, որ շարժումը քոնն է, ու զգացողությունը «անտեսում» է
        \item Սա ներքինն արտաքինից տարբերելու լավ օրինակ է
      \end{itemize}}
    \end{column}
    \begin{column}{0.44\textwidth}
      \onslide<3->{%
      \begin{card}[colback=lavender!40!cardwhite]
        \emoji[1.3em]{feather}\ {\blackfont Բլեյքմորի փորձը}\par\smallskip
        Հալյուցինացիաներ ունեցող պացիենտներն ինքնախուտուտի նկատմամբ \textbf{շատ ավելի զգայուն} էին, քան մյուսները
      \end{card}}
    \end{column}
  \end{columns}
\end{frame}

\begin{frame}{Բարձրահասակ ընկերոջ դեպքը}
  \begin{columns}[c]
    \begin{column}{0.56\textwidth}
      \begin{itemize}
        \item Զգուշացում. եթե կարողանում ես ինքդ քեզ խուտուտ տալ, դա \textbf{չի նշանակում}, որ փսիխոզ ունես։ Մարդիկ շատ տարբեր են
        \item<2-> Բըրնեթի կնոջ ուսանողական տարիների ընկերը կարողանում էր, ու ոչ մի հոգեկան խնդիր չուներ։ Բայց շատ բարձրահասակ էր
        \item<3-> Բըրնեթի «վարկածը». ազդանշանն այնքան երկար է գնում ուղեղ, որ ուղեղը մոռանում է՝ որտեղից է եկել
      \end{itemize}
    \end{column}
    \begin{column}{0.38\textwidth}
      \onslide<4->{%
      \begin{bookquote}
        Սա բացարձակապես անհնար է։ Ուսանող էի, հանկարծակի հարցրին, ու այն ժամանակ ես շատ ավելի ամբարտավան էի. նախընտրում էի խենթ վարկած հորինել, քան խոստովանել, որ չգիտեմ։
      \end{bookquote}}
    \end{column}
  \end{columns}
\end{frame}

\begin{frame}{Ինչպես է առաջանում հալյուցինացիան}
  \centering
  % How the brain tells inside from outside, after Allen et al. (2008), as the book describes it.
% Slide 1: the normal network. Slide 2: a hallucination (internal signal too strong, monitoring links weakened).
\begin{tikzpicture}[nd/.style={flow=3.2cm, font=\small, inner sep=3pt}]
  \node[nd, fill=butter!50!cardwhite] (th) at (-4.8,0) {\textbf{Թալամուս}\\ «հում» տվյալներ զգայարաններից};
  \node[nd, fill=lavender!45!cardwhite, text width=6.4cm] (mon) at (0,2.15)
    {\textbf{Նախաճակատային, նախաշարժիչ կեղև, ուղեղիկ}\\ «Սա \textbf{ես} եմ անում»};
  \node[nd, fill=mint!45!cardwhite, text width=4.4cm] (em) at (0,-2.0) {\textbf{Հիպոկամպ, նշաձև մարմին}\\ հիշողություն, զգացմունք};
  \only<1>{\node[nd, fill=cardwhite, text width=3.4cm] (sc) at (0,0)
    {\textbf{Զգայական կեղև}\\ գծեր, գույներ $\rightarrow$\\ «ավտոբուս է»};}
  \only<2->{\node[nd, fill=amOrange!35!cardwhite, draw=amOrange!80!black, text width=3.4cm] (sc) at (0,0)
    {\textbf{Զգայական կեղև}\\ ներքինը \textbf{ուժեղացել} է};}
  \only<1>{\node[nd, fill=cardwhite, text width=2.6cm] (out) at (4.8,0) {\emoji{eye}\emoji{ear}\\ ընկալում};}
  \only<2->{\node[nd, fill=brainpink!45!cardwhite, text width=2.6cm] (out) at (4.8,0) {\emoji{speakinghead}\\ միտքը թվում է \textbf{դրսից}};}
  \draw[arr] (th) -- (sc);
  \draw[arr] (sc) -- (out);
  \draw[arr, <->] (em) -- (sc);
  \only<1>{\draw[arr, <->] (mon) -- (sc);}
  \only<2->{\draw[arr=amRed, <->, dashed] (mon) -- (sc);
    \node[font=\blackfont\scriptsize, text=amRed, anchor=west, inner sep=1pt] at (0.15,1.1) {կապը թուլացել է};
    \node[font=\blackfont\scriptsize, text=amRed, anchor=west] at (em.east) {զգացմունքը չի մեղմվում};}
\end{tikzpicture}

\end{frame}

\begin{frame}{«Սա հիմարություն էր»}
  \begin{columns}[c]
    \begin{column}{0.48\textwidth}
      {\small Թանկ թեյի սերվիզ ես գնել ու այն տվել ես փոքրիկիդ, որ տանի խանութից դուրս...}\par\medskip
      \centering
      \begin{tikzpicture}
        \node at (0,0) {\mascot[1.9cm]{confused}};
        \node[draw=plum, dashed, line width=0.8pt, rounded corners=3mm, fill=cardwhite, text=plum, align=center,
              font=\handfont\large, text width=2.9cm, inner sep=4pt, visible on=<1>] at (2.9,0.9) {Սա հիմարություն էր};
        \filldraw[fill=cardwhite, draw=plum, line width=0.6pt, visible on=<1>] (1.05,0.45) circle (0.07) (1.3,0.6) circle (0.1);
        \node[say, text width=2.9cm, font=\handfont\large, callout absolute pointer={(0.8,0.4)}, fill=brainpink!20!cardwhite, visible on=<2->] at (2.9,0.9) {«Սա հիմարություն էր»};
      \end{tikzpicture}
    \end{column}
    \begin{column}{0.48\textwidth}
      \begin{itemize}
        \item Սովորաբար սա \textbf{ներքին} միտք է
        \item<2-> Եթե ուղեղը չճանաչի, որ միտքն իրենն է, լեզվի կենտրոններն այն կընկալեն որպես \textbf{խոսք}
        \item<3-> Իսկ եթե նշաձև մարմինը չմեղմի զգացմունքը, ձայնը կթվա շատ \textbf{քննադատող}
        \item<4-> Նույնը կարող է լինել \textbf{բոլոր} զգայարանների հետ
      \end{itemize}
    \end{column}
  \end{columns}
\end{frame}

\begin{frame}{Զառանցանքներ}
  \begin{columns}[T]
    \begin{column}{0.47\textwidth}
      \begin{card}[colback=butter!45!cream]
        \emoji[1.3em]{crown}\ {\blackfont Մեծամտության} (grandiose)\par\smallskip
        Կոշիկի խանութի կես դրույքով աշխատողը վստահ է, որ աշխարհի լավագույն գործարարն է
      \end{card}
    \end{column}
    \begin{column}{0.47\textwidth}
      \onslide<2->{%
      \begin{card}[colback=lavender!40!cardwhite]
        \emoji[1.3em]{detective}\ {\blackfont Հետապնդման} (persecutory)\par\smallskip
        Ավելի հաճախ է. «բոլորը մասնակցում են ինձ առևանգելու գաղտնի դավադրությանը»
      \end{card}}
    \end{column}
  \end{columns}
  \vspace{3mm}
  \onslide<3->{%
  \begin{card}[colback=brainpink!25!cardwhite]
    Զառանցանքները շատ \textbf{համառ} են։ Ավելի հեշտ է համոզել, որ ձայներն իրական չեն, քան որ ոչ բոլորն են քո դեմ դավադրություն անում
  \end{card}}
\end{frame}

\begin{frame}{Աշխարհի մոդելը. ավտոբուսը տան մոտ}
  \begin{columns}[c]
    \begin{column}{0.46\textwidth}
      {\small Ուղեղը պահում է աշխարհի \textbf{մոդելը}, որ ամեն պահ չզարմանանք։ Ավտոբուսը կանգառում կանգնում է. նորմալ է}\par\medskip
      \begin{chatbox}
        Ավտոբուսը կանգնել է \textbf{քո տան} դիմաց ու չի շարժվում։ Ի՞նչ բացատրություններ են գալիս մտքիդ։
      \end{chatbox}
    \end{column}
    \begin{column}{0.5\textwidth}
      \onslide<2->{%
      \begin{popup}{ուղեղ.exe}
        \small
        \begin{itemize}\setlength{\itemsep}{0.3mm}
          \item Վարորդն ինձ լրտեսու՞մ է
          \item Ինչ-որ մեկն ինձ ավտոբուս է նվիրե՞լ
          \item Տունս դարձե՞լ է ավտոբուսների կայան
        \end{itemize}
      \end{popup}}
      \onslide<3->{%
      \begin{card}[colback=mint!35!cardwhite]
        Ուղեղը դրանք մերժում է որպես անհավանական։ Պարզվում է՝ ավտոբուսը \textbf{փչացել է}
      \end{card}}
    \end{column}
  \end{columns}
\end{frame}

\begin{frame}{Կապգրայի զառանցանքը}
  \relax{\small Մարդը վստահ է, որ մտերիմը (կինը, ծնողը, ընկերը, նույնիսկ ընտանի կենդանին) փոխարինվել է \textbf{նույն տեսքով խաբեբայով}}\par\medskip
  \centering
  % Capgras delusion as the book explains it: the expected feelings are missing, so the brain resolves the mismatch
% with "an identical impostor". Steps reveal on slides 1-5.
\begin{tikzpicture}[st/.style={flow=2.45cm, font=\small, inner sep=3pt, minimum height=1.55cm}]
  \node[st, fill=butter!50!cardwhite] (a) at (0,0) {\emoji{eyes}\\ «Սա իմ կինն է»};
  \node[st, fill=lavender!40!cardwhite, visible on=<2->] (b) at (2.85,0) {\emoji{heart}\\ պիտի սեր, ջերմություն զգամ};
  \node[st, fill=peach!45!cardwhite, visible on=<3->] (c) at (5.7,0) {\emoji{neutral}\\ ոչինչ չեմ զգում};
  \node[st, fill=brainpink!35!cardwhite, visible on=<4->] (d) at (8.55,0) {\emoji{question}\\ ինչո՞ւ};
  \node[st, fill=amRed!20!cardwhite, visible on=<5->] (e) at (11.4,0) {\emoji{detective}\\ «Նա խաբեբա է՝ նույն տեսքով»};
  \draw[arr, visible on=<2->] (a) -- (b);
  \draw[arr, visible on=<3->] (b) -- (c);
  \draw[arr, visible on=<4->] (c) -- (d);
  \draw[arr, visible on=<5->] (d) -- (e);
  \node[font=\scriptsize, text=plum!75!cream, align=center, visible on=<3->] at (5.7,-1.15) {ճակատային բլթի վնաս};
\end{tikzpicture}
\par\medskip
  \raggedright
  \onslide<6->{%
  \begin{card}[colback=lavender!35!cardwhite]
    Ապացույցները չեն օգնում. «Դա իսկապես ես եմ» լսելուց հետո զգացմունքների բացակայությունն ավելի է աչքի ընկնում, ու «խաբեբայի» բացատրությունն ավելի համոզիչ է թվում
  \end{card}}
\end{frame}

\begin{frame}{Ինչպես են ծնվում զառանցանքները}
  \centering
  \begin{tikzpicture}[st/.style={flow=2.8cm, font=\small}]
    \node[st, fill=butter!50!cardwhite] (a) at (0,0) {սպասում եմ\\ մի բան};
    \node[st, fill=lavender!40!cardwhite] (b) at (3.4,0) {տեսնում եմ\\ ուրիշ բան};
    \node[st, fill=peach!45!cardwhite] (c) at (6.8,0) {չի\\ համընկնում};
    \node[st, fill=brainpink!40!cardwhite] (d) at (10.2,0) {պետք է\\ բացատրություն};
    \draw[arr] (a) -- (b); \draw[arr] (b) -- (c); \draw[arr] (c) -- (d);
  \end{tikzpicture}\par\smallskip
  \raggedright
  \begin{columns}[T]
    \begin{column}{0.47\textwidth}
      \onslide<2->{%
      \begin{card}[colback=peach!35!cardwhite]
        \emoji[1.3em]{fearful}\ {\blackfont Անհանգստություն}\par\smallskip
        Վտանգի համակարգը միացած է առանց պատճառի, ու ուղեղը «գտնում է» վտանգը. խանութում մրմնջացող կինը դավադրության մաս է
      \end{card}}
    \end{column}
    \begin{column}{0.47\textwidth}
      \onslide<3->{%
      \begin{card}[colback=lavender!35!cardwhite]
        \emoji[1.3em]{pensive}\ {\blackfont Դեպրեսիա}\par\smallskip
        Անբացատրելի վատ տրամադրություն։ Մեկը վեր կացավ, երբ նստեցիր. «Ուրեմն ինձ ատում են»
      \end{card}}
    \end{column}
  \end{columns}
\end{frame}

\begin{frame}{Ինչու են զառանցանքներն այդքան համառ}
  \begin{columns}[c]
    \begin{column}{0.54\textwidth}
      \begin{itemize}
        \item Ինչը չի համապատասխանում մեր մոդելին, ուղեղը սովորաբար անտեսում է
        \item<2-> Զառանցանքի ժամանակ հակառակ ապացույցները ճնշվում են \textbf{ավելի ուժեղ}
        \item<3-> Գուցե մեղավոր է \textbf{գլուտամատի} (գրգռող նեյրոմեդիատոր) ավելցուկը. անմեղ բաները թվում են կարևոր
        \item<4-> Ավելորդ ակտիվությունը թուլացնում է պլաստիկությունը, ու զառանցանքը «սառչում» է
      \end{itemize}
    \end{column}
    \begin{column}{0.4\textwidth}
      \onslide<2->{%
      \begin{card}[colback=mint!35!cardwhite]
        \emoji[1.3em]{alien}\ Չես հավատում այլմոլորակայիններին, ու «առևանգվածների» պատմությունները քեզ չեն համոզում։\par\smallskip
        Բայց եթե քեզ առևանգեն ու մանրակրկիտ հետազոտեն, կարծիքդ կփոխվի։ Զառանցանքի ժամանակ նույնիսկ դա չի օգնում
      \end{card}}
    \end{column}
  \end{columns}
\end{frame}

\begin{frame}{Ոչ ամեն տարօրինակ համոզմունք է զառանցանք}
  \begin{columns}[c]
    \begin{column}{0.54\textwidth}
      \begin{itemize}
        \item Միլիոնավոր մարդիկ հավատում են, որ Երկիրը 6000 տարեկան է, ու դինոզավրեր չեն եղել
        \item Ոմանց հետ «խոսում» են մահացած հարազատները։ Հիվանդությու՞ն է, սու՞գ, հոգևոր փորձառությու՞ն
        \item<2-> Ուղեղը «իրականը» որոշում է \textbf{փորձով} ու միջավայրով
      \end{itemize}
    \end{column}
    \begin{column}{0.4\textwidth}
      \onslide<3->{%
      \begin{card}[colback=butter!45!cream]
        Զառանցանք է միայն այն համոզմունքը, որը \textbf{չի համապատասխանում} մարդու հավատալիքներին։\par\smallskip
        Բարեպաշտ քարոզիչը լսում է Աստծո ձայնը. զառանցանք չէ։ Ագնոստիկ հաշվապահ ուսանողուհին՝ հավանաբար այո
      \end{card}}
    \end{column}
  \end{columns}
\end{frame}

\begin{frame}{Իսկ եթե ամբողջ գիրքը հալյուցինացիա էր}
  \begin{columns}[c]
    \begin{column}{0.58\textwidth}
      \begin{itemize}
        \item Մեր ընկալումը հաշվարկներ են ու երբեմն պարզապես \textbf{գուշակություն}
        \item «Նորմալը» ավելի շատ համաձայնություն է, քան փաստ
      \end{itemize}
      \onslide<2->{%
      \begin{bookquote}
        Գուցե ոչինչ իրական չէ՞։ Գուցե այս ամբողջ գիրքը հալյուցինացիա՞ էր։ Հուսով եմ՝ ոչ, այլապես բավական շատ ժամանակ ու ջանք եմ վատնել։
      \end{bookquote}}
    \end{column}
    \begin{column}{0.36\textwidth}
      \centering\mascot[2.8cm]{confused}
    \end{column}
  \end{columns}
\end{frame}
